\subsection{TRACE OF A MATRIX}

Let C be a commutative ring; for every square matrix \(X = (x_{ij})\) over C corresponding to the finite indexing set I, the trace of X is the element

\[
\operatorname{Tr} (X) = \sum_ {i \in I} x _ {i i}.\tag{49}
\]

Let E be a C-module admitting a finite basis \((e_{i})_{i\in I}\) ; for every endomorphism u of E,

\[
\operatorname{Tr} (u) = \operatorname{Tr} (M (u))\tag{50}
\]

\(M(u)\) being the matrix u with respect to the basis \((e_{i})\) ; this follows immediately from § 4, no. 3, formula (17), when this formula is applied to the endomorphism \(x \mapsto \langle x, e_{i}^{*} \rangle e_{j}\) (where \(e_{i}^{*}\) ) is the dual basis of \((e_{i})\) ); from this we pass to the general case by linearity. Formula (49) shows that

\[
\operatorname{Tr} (u) = \sum_ {i} \left\langle u (e _ {i}), e _ {i} ^ {*} \right\rangle\tag{51}
\]

for every basis \((e_i)\) of E (cf.~§4, no. 3, formula (17)).

If X is a matrix of type \((m, n)\) over C and Y a matrix of type \((n, m)\) over C, then

\[
\operatorname{Tr} (X Y) = \operatorname{Tr} (Y X)\tag{52}
\]

as follows from the above and Proposition 3 of § 4, no. 3; (52) can also be obtained directly, for if \(X = (x_{ij})\) , \(Y = (y_{ji})\) ( \(1 \leqslant i \leqslant m\) , \(1 \leqslant j \leqslant n\) ), then

\[
\operatorname{Tr} (X Y) = \sum_ {i, j} x _ {i j} y _ {j i}\tag{53}
\]

by (49). The latter formula proves moreover:

PROPOSITION 7. Let C be a commutative ring and for every matrix \(P \in \mathbf{M}_n(\mathbf{C})\) let \(f_P\) be the linear form \(X \mapsto \operatorname{Tr}(PX)\) on \(\mathbf{M}_n(\mathbf{C})\) ; the mapping \(P \mapsto f_P\) is a C-linear bijection of \(\mathbf{M}_n(\mathbf{C})\) onto its dual.

PROPOSITION 8. If \(g\) is a linear form on the C-module \(\mathbf{M}_n(\mathbf{C})\) such that \(g(XY) = g(YX)\) for all matrices \(X, Y\) in \(\mathbf{M}_n(\mathbf{C})\) , there exists one and only one scalar \(c \in \mathbf{C}\) such that \(g(X) = c\) . Tr(X) for every matrix \(X \in \mathbf{M}_n(\mathbf{C})\) .

Since the proposition is trivial for \(n = 1\) , attention may be confined to the case where \(n \geqslant 2\) . Taking \(X = E_{ij}\) , \(Y = E_{jk}\) with \(i \neq k\) , we obtain \(g(E_{ik}) = 0\) ; then taking \(X = E_{ij}\) , \(Y = E_{ji}\) with \(i \neq j\) , we find \(g(E_{ii}) = g(E_{jj})\) ; the proposition follows immediately since the \(E_{ij}\) form a basis of \(\mathbf{M}_n(\mathbf{C})\) .

